\documentclass[]{article}

% Imports the catppuccin theme, using the mocha flavor,
% from the directory above. Actual implementation
% wouldn't need the import package unless the theme
% and the document are in different directories.
\usepackage{import}
\usepackage{xcolor}
\usepackage{cancel}
\usepackage{mathtools}

% For permutations and combinations
\newcommand\Myperm[2][^n]{\prescript{#1\mkern-2.5mu}{}P_{#2}}
\newcommand\Mycomb[2]{\prescript{#1\mkern-0.5mu}{}C_{#2}}

% Colors
\definecolor{yorhabg}{HTML}{131314}
\definecolor{yorhafg}{HTML}{C9C7CD}
\definecolor{yorhagrid}{HTML}{B5AF9C}
\definecolor{mred}{HTML}{D67069}
\definecolor{mblue}{HTML}{6887A1}

\pagecolor{yorhabg}
\color{yorhafg}

\usepackage{preamble}

% Removes padding above title
\usepackage{titling}
\setlength{\droptitle}{-10em}

% Font package
\usepackage[T1]{fontenc}

\usepackage{fouriernc}

\usepackage{sectsty}
\usepackage{graphicx}
\usepackage{amsmath}
\usepackage{amsfonts}
\usepackage{amssymb}
\usepackage[skins, most]{tcolorbox}

\DeclareMathOperator{\sgn}{sgn}

\usepackage{tikz}
\usepackage{eso-pic}
\usetikzlibrary{calc,shadows.blur}
\usetikzlibrary{angles, quotes}
\usetikzlibrary{3d}

% Margins
\topmargin=0in
\evensidemargin=0in
\oddsidemargin=0in
\textwidth=6.5in
\textheight=9.0in
\headsep=0.25in

\AtBeginEnvironment{tcolorbox}{\small}

\newtcolorbox{imp}{enhanced,arc=0mm,colback=yorhabg,colframe=mred,leftrule=10mm,coltext=yorhafg,%
overlay={\node[anchor=west,outer sep=2pt] at (frame.west) {\includegraphics[width=6mm]{images/imageb.png}}; }}

\newtcolorbox{shortcut}{enhanced,arc=0mm,colback=yorhabg,colframe=mred,leftrule=10mm,coltext=yorhafg, coltitle=yorhabg, title=\texttt{Shortcut.}, 
overlay={\node[anchor=west,outer sep=2pt] at (frame.west) {\includegraphics[width=6mm]{images/imageb.png}}; }}

\newtcolorbox{question}{
    enhanced, 
    colback=yorhabg,
    colframe=mblue,
    coltext=yorhafg,
    coltitle=yorhabg,
    attach boxed title to top left={yshift*=-\tcboxedtitleheight}, 
    title=\texttt{Question.},
    boxed title size=title,
    boxed title style={%
        rounded corners=northeast, 
        rounded corners=northwest, 
        colback=tcbcolframe, 
        boxrule=0pt,
    },
    underlay boxed title={%
        \path[fill=tcbcolframe] (title.south west)--(title.south east) 
            to[out=0, in=180] ([xshift=5mm]title.east)--
            (title.center-|frame.east)
            [rounded corners=5pt] |- 
            (frame.north) -| cycle; 
    },
}

\newcommand\bb[1]{\textcolor{yorhafg}{\textbf{#1}}}

\title{\textbf{CSCA67 - Assignment \#1}}
\author{ Satyajit Datta }
\date{\today}

\begin{document}

\maketitle    

\section*{1. Logical Equivalence}
For each of the following pairs of expressions, either prove that the two expressions are equivalent or prove that
they are not. (Clearly state what you are proving!) Do not use truth tables.
\vspace{1em}
\hrule
\vspace{1em}

\begin{question}
    $(a \rightarrow b) \land (b \rightarrow c)$ and $a \rightarrow c$
\end{question}
\begin{center}
        \text{We are proving that the two expressions are not equivalent.} \\
        \text{Consider the case where } a = \text{True}, b = \text{False}, c = \text{True}. \\
\end{center}
\begin{align*}
    & (a \rightarrow b) \land (b \rightarrow c) & & a \rightarrow c\\
    \Rightarrow\; & (T \rightarrow F) \land (F \rightarrow T) & \Rightarrow\; & T \rightarrow T\\
    \Rightarrow\; & F \land T & \Rightarrow\; & \bb{T}\\
    \Rightarrow\; & \bb{F} && \\
\end{align*}
\begin{center}
    $\therefore$ The two expressions are not equivalent. $\blacksquare$
\end{center}

\begin{question}
    $a \land (a \rightarrow b)$ and $a \rightarrow b$
\end{question}
\begin{center}
    \text{We are proving that the two expressions are not equivalent.} \\
    \text{Consider the case where} a = \text{False}, b = \text{True}
\end{center}
\begin{align*}
    & a \land (a \rightarrow b) &  & a \rightarrow b \\
    \Rightarrow\; & F \land (F \rightarrow T)  & \Rightarrow\; & F \rightarrow T\\
    \Rightarrow\; & F \land T & \Rightarrow\; \bb{T} \\
    \Rightarrow\; & \bb{F}
\end{align*}
\begin{center}
    $\therefore$ The two expressions are not equivalent. $\blacksquare$
\end{center}

\begin{question}
    $(a \rightarrow b) \land (a \rightarrow c)$ and $a \rightarrow (b \land c)$
\end{question}
\begin{center}
    \text{We are proving that the two expressions \bb{are} equivalent.}
\end{center}
\begin{align*}
    & (a \rightarrow b) \land (a \rightarrow c) & & a \rightarrow (b \land c) \\
    \Rightarrow\; & (\neg a \lor b) \land (\neg a \lor c) \quad \text{(Conditional Law)} & \Rightarrow\; & \mathbf{\neg a \lor (b \land c)} \quad \text{(Conditional Law)} \\
    \Rightarrow\; & \mathbf{\neg a \lor (b \land c)} \quad \text{(Distributive Law)} & & \\
\end{align*}
\begin{center}
    $\therefore$ The two expressions are equivalent. $\blacksquare$
\end{center}

\begin{question}
    $(a \rightarrow c) \lor (b \rightarrow c)$ and $(a \lor b) \rightarrow c$
\end{question}
\begin{center}
    \text{We are proving that the two expressions \bb{are} equivalent.}
\end{center}
\begin{align*}
    & (a \rightarrow c) \lor (b \rightarrow c) & & (a \lor b) \rightarrow c \\
    \Rightarrow\; & (\neg a \lor c) \land (\neg b \lor c) \quad \text{(Conditional Law)} & \Rightarrow\; & \neg(a \lor b) \lor c \quad \text{(Conditional Law)} \\
    \Rightarrow\; & \mathbf{(\neg a \land \neg b) \lor c} \quad \text{(Distributive Law)} & \Rightarrow\; & \mathbf{(\neg a \land \neg b) \lor c} \quad \text{(De Morgan's Theorem)} \\
\end{align*}
\begin{center}
    $\therefore$ The two expressions are equivalent. $\blacksquare$
\end{center}

\begin{question}
    $a \iff b$ and $(a \land b) \lor (\neg a \land \neg b)$
\end{question}
\begin{center}
    \text{We are proving that the two expressions \bb{are} equivalent.}
\end{center}
\begin{align*}
    & a \iff b & & \mathbf{(a \land b) \lor (\neg a \land \neg b)} \\
    \Rightarrow\; & (a \rightarrow b) \land (b \rightarrow a) \quad \text{(Biconditional Law)} \\
    \Rightarrow\; & (\neg a \lor b ) \land (\neg b \lor a) \quad \text{(Conditional Law)} \\
    \Rightarrow\; & (\neg a \land \neg b) \lor (\neg a \land a) \lor (b \land \neg b) \lor (b \land a) \quad \text{(Distributive Law)} \\ 
    \Rightarrow\; & \mathbf{(\neg a \land \neg b) \lor (a \land b)} \quad \text{(Negation Law)}
\end{align*}
\begin{center}
    $\therefore$ The two expressions are equivalent. $\blacksquare$
\end{center}

\begin{question}
    $a \rightarrow (b \rightarrow (c \rightarrow d))$ and $(a \land b \land c) \rightarrow d$
\end{question}
\begin{center}
    \text{We are proving that the two expressions \bb{are} equivalent.}
\end{center}
\begin{align*}
    & a \rightarrow (b \rightarrow (c \rightarrow d)) & & \mathbf{(a \land b) \lor (\neg a \land \neg b)} \\
    \Rightarrow\; & \neg a \lor (\neg b \lor (\neg c \lor d)) \quad \text{(Conditional Law)} \\
    \Rightarrow\; & (\neg a \lor \neg b \lor \neg c) \lor d \\
    \Rightarrow\; & \neg(a \land b \land c) \lor d \quad \text{(De Morgan's)} \\
    \Rightarrow\; & \mathbf{(a \land b \land c) \rightarrow d} \quad \text{(Conditional Law)} \\
\end{align*}
\begin{center}
    $\therefore$ The two expressions are equivalent. $\blacksquare$
\end{center}

\begin{question}
    $(a \rightarrow b) \lor (b \rightarrow a)$ and $a \iff b$
\end{question}
\begin{center}
        \text{We are proving that the two expressions are not equivalent.} \\
        \text{Consider the case where } a = \text{True}, b = \text{False}. \\
\end{center}
\begin{align*}
    & (a \rightarrow b) \land (b \rightarrow a) & & a \iff b\\
    \Rightarrow\; & (T \rightarrow F) \lor (F \rightarrow T) & \Rightarrow\; & (a \rightarrow b) \land (b \rightarrow a) \quad \text{(Biconditional Law)}\\
    \Rightarrow\; & F \lor T & \Rightarrow\; &  (T \rightarrow F) \land (F \rightarrow T)\\
    \Rightarrow\; & \mathbf{T} & \Rightarrow\; & F \land T \\
    & &  \Rightarrow\; & \mathbf{F} \\
\end{align*}
\begin{center}
    $\therefore$ The two expressions are not equivalent. $\blacksquare$
\end{center}

\begin{question}
    $a \iff b$ and $\neg a \iff \neg b$
\end{question}
\begin{center}
    \text{We are proving that the two expressions \bb{are} equivalent.}
\end{center}
\begin{align*}
    & a \iff b & & \neg a \iff \neg b \\
    \Rightarrow\; & (a \rightarrow b) \land (b \rightarrow a) \quad \text{(Biconditional Law)} & \Rightarrow & (\neg a \rightarrow \neg b) \land (\neg b \rightarrow \neg a) \quad \text{(Biconditional Law)}\\
    \Rightarrow\; & (\neg a \lor b) \land (\neg b \lor a) \quad \text{(Conditional Law)} & \Rightarrow & (\neg\neg a \lor \neg b) \land (\neg\neg b \lor \neg a) \quad \text{(Conditional Law)}\\
    & & \Rightarrow & (a \lor \neg b) \land (b \lor \neg a) \quad \text{(Double Negation law)} \\
\end{align*}
\begin{center}
    $\therefore$ The two expressions are equivalent. $\blacksquare$
\end{center}

\begin{question}
    $(a \land b) \rightarrow (c \land d)$ and $((a \rightarrow c) \land (a \rightarrow d)) \land ((b \rightarrow c) \land (b \rightarrow d))$
\end{question}
\begin{center}
        \text{We are proving that the two expressions are not equivalent.} \\
        \text{Consider the case where } a = \text{True}, b = \text{False}, c = \text{False}, d = \text{False}.\\
\end{center}
\begin{align*}
    & (a \land b) \rightarrow (c \land d) & & ((a \rightarrow c) \land (a \rightarrow d)) \land ((b \rightarrow c) \land (b \rightarrow d))\\
    \Rightarrow\; & (T \land F) \rightarrow (F \land F) & \Rightarrow\; & ((T \rightarrow F) \land (T \rightarrow F)) \land ((F \rightarrow F) \land (F \rightarrow F)) \\
    \Rightarrow\; & F \lor T & \Rightarrow\; &  (F \land F) \land (T \land T)\\
    \Rightarrow\; & \mathbf{T} & \Rightarrow\; & F \land F \\
    & &  \Rightarrow\; & \mathbf{F} \\
\end{align*}
\begin{center}
    $\therefore$ The two expressions are not equivalent. $\blacksquare$
\end{center}

\begin{question}
        $(a \lor b) \rightarrow (c \land d)$ and $((a \rightarrow c) \land (a \rightarrow d)) \land ((b \rightarrow c) \land (b \rightarrow d))$
\end{question}
\begin{center}
        \text{We are proving that the two expressions are equivalent.} \\
\end{center}
\begin{align*}
    & (a \lor b) \rightarrow (c \land d) & & ((a \rightarrow c) \land (a \rightarrow d)) \land ((b \rightarrow c) \land (b \rightarrow d))\\
    \Rightarrow\; & \neg(a \lor b) \lor (c \land d) \quad \text{(Conditional Law)} & \Rightarrow\; & ((\neg a \lor c) \land (\neg a \lor d)) \land ((\neg b\rightarrow c) \land (\neg b \lor d)) \quad \text{(Conditional Law)}\\
    &  & \Rightarrow\; &  (\neg a \lor (c \land d))\land (\neg b \lor (c \land d)) \quad \text{(Distributive Law)}\\
    & & \Rightarrow\; & (\neg a \land \neg b ) \lor (c \land d) \quad \text{(Distributive Law)} \\
    & &  \Rightarrow\; & \neg(a \lor b) \lor (c \land d) \quad \text{(De Morgan's law)} \\
\end{align*}
\begin{center}
    $\therefore$ The two expressions are equivalent. $\blacksquare$
\end{center}

\section*{2. Logical Equivalence}
Determine whether each statement below is a tautology, a contradiction, or neither. Prove your result. Do not
use truth tables.

\begin{center}
    Note that for this section, I will refer to tautologies with $\top$ and contradictions with $\bot$.
\end{center}
\vspace{0.1in}
\hrule
\vspace{0.1in}
\begin{question}
    $(a \rightarrow b) \lor (b\ \rightarrow a)$
\end{question}
\begin{align*}
    & (a \rightarrow b) \lor (b \rightarrow a) \\
    \Rightarrow\; & (\neg a \lor b) \lor (\neg b \lor a) \quad \text{(Conditional Law)} \\
    \Rightarrow\; & (\neg a \lor a) \lor (\neg b \lor b) \quad \\
    \Rightarrow\;&  (t)  \lor (t) \quad \text{(Tautology Law)} \\
    \Rightarrow\; & \mathbf{t} \quad \text{(Tautology Law)}
\end{align*}
\begin{center}
    $\therefore$ The statement is a tautology. $\blacksquare$
\end{center}

\begin{question}
    $(((a \rightarrow b) \land (b \rightarrow c)) \land a) \land \neg c$
\end{question}
\begin{align*}
    & (((a \rightarrow b) \land (b \rightarrow c)) \land a) \land \neg c \\
    \Rightarrow\; & (((\neg a \lor b) \land (\neg b \lor c)) \land a) \land \neg c \quad \text{(Conditional Law)} \\
    \Rightarrow\; & ((\neg a \land a) \lor (a\land b) \land (\neg b \lor c)) \land \neg c \quad \text{(Distributive Law)} \\
    \Rightarrow\; & (\bot \lor (a\land b) \land (\neg b \lor c)) \land \neg c \quad \text{(Contradiction Law)} \\
    \Rightarrow\; & ((a\land b) \land (\neg b \lor c)) \land \neg c \quad \text{(Contradiction Law)} \\
    \Rightarrow\; & ((a \land \neg b \land b) \lor (a \land b \land c)) \land \neg c \quad \text{(Distributive Law)} \\
    \Rightarrow\; & ((a \land \bot) \lor (a\land b\land c)) \land \neg c \quad \text{(Contradiction Law)} \\
    \Rightarrow\; & (\bot \lor (a\land b\land c)) \land \neg c \quad \text{(Contradiction Law)} \\
    \Rightarrow\; & (a\land b\land c) \land \neg c \quad \text{(Contradiction Law)} \\
    \Rightarrow & a \land b \land c \land \neg c \quad \\
    \Rightarrow\; & a \land b \land \bot \quad \text{(Contradiction Law)} \\
    \Rightarrow\; & \bot \quad \text{(Contradiction Law)}
\end{align*}
\begin{center}
    $\therefore$ The statement is a contradiction. $\blacksquare$
\end{center}
\vspace{1.7in}
\begin{question}
    $(((a \rightarrow b) \land (b \rightarrow c)) \land a) \land  c$
\end{question}
\begin{align*}
    & (((a \rightarrow b) \land (b \rightarrow c)) \land a) \land c \\
    \Rightarrow\; & (((\neg a \lor b) \land (\neg b \lor c)) \land a) \land c \quad \text{(Conditional Law)} \\
    \Rightarrow\; & ((\neg a \land a) \lor (a\land b) \land (\neg b \lor c)) \land c \quad \text{(Distributive Law)} \\
    \Rightarrow\; & (\bot \lor (a\land b) \land (\neg b \lor c)) \land c \quad \text{(Contradiction Law)} \\
    \Rightarrow\; & ((a\land b) \land (\neg b \lor c)) \land c \quad \text{(Contradiction Law)} \\
    \Rightarrow\; & ((a \land \neg b \land b) \lor (a \land b \land c)) \land  c \quad \text{(Distributive Law)} \\
    \Rightarrow\; & ((a \land \bot) \lor (a\land b\land c)) \land  c \quad \text{(Contradiction Law)} \\
    \Rightarrow\; & (\bot \lor (a\land b\land c)) \land c \quad \text{(Contradiction Law)} \\
    \Rightarrow\; & (a\land b\land c) \land c \quad \text{(Contradiction Law)} \\
    \Rightarrow\; & a \land b \land c \land c \quad \\
    \Rightarrow\; & a \land b \land c \quad \text{(Idempotent Law)}
\end{align*}
\begin{center}
    $\therefore$ The statement is a neither a tautology nor a contradiction. $\blacksquare$
\end{center}

\begin{question}
    $a \rightarrow \neg a$
\end{question}
\begin{align*}
    & a \rightarrow \neg a \\
    \Rightarrow\; & \neg a \lor \neg a \quad \text{(Conditional Law)} \\
    \Rightarrow\; & \neg a \quad \text{(Idempotent Law)} \\
\end{align*}
\begin{center}
    $\therefore$ Therefore, the statement is neither a tautology nor a contradiction. $\blacksquare$
\end{center}

\begin{question}
    $(a \land (a \rightarrow b)) \rightarrow b$
\end{question}
\begin{align*}
    \Rightarrow\; & (a \land (a \rightarrow b)) \rightarrow b \\
    \Rightarrow\; & \neg(a \land (a \rightarrow b)) \lor b \quad \text{(Conditional Law)} \\
    \Rightarrow\; & \neg(a \land (\neg a \lor b)) \lor b \quad \text{(Conditional Law)} \\
    \Rightarrow\; & \neg((a \land \neg a) \lor (a \land b)) \lor b \quad \text{(Distributive Law)} \\
    \Rightarrow\; & \neg(\top \lor (a\land b)) \lor b \quad \text{(Tautology Law)}\\
    \Rightarrow\; &  \neg(a \land b) \lor b \quad \text{(Tautology Law)} \\
    \Rightarrow\; & \neg a \lor \neg b \lor b \quad \text{(De Morgan's Law)} \\
    \Rightarrow\; & \neg a \lor \top \quad \text{(Tautology Law)} \\
    \Rightarrow\; & \top \quad \text{(Tautology Law)}
\end{align*}
\begin{center}
    $\therefore$ The statement is a tautology. $\blacksquare$
\end{center}
\section*{3. Deductive Reasoning}
For each of the following arguments:

\quad • If the argument is valid, you need to both:

\quad \quad 1. Prove the argument is valid by using a truth table. The truth table must be in the format we used in class, and the variables must appear in alphabetical order.

\quad \quad 2. Prove the argument is valid by using Rules of Inference.

\quad • If the argument is not valid, provide an assignment of truth values to variables and demonstrate that it proves
that the argument is not valid.
\vspace{0.1in}
\hrule
\vspace{0.1in}

JOISADJFIOJASF KILL ME NOW
\end{document}